%=========== ΛΥΣΗ - 69ο ΘΕΜΑ Δ ===========
\begin{thema}{Δ}
\begin{erwthma}
\item Από την εκφώνηση γνωρίζουμε ότι η εφαπτομένη της $C_f$ στο
$A(1,3)$ έχει κλίση $2$, άρα
\[
f'(1)=2.
\]
Επίσης, η εφαπτομένη της $C_f$ στο $B(4,1)$ έχει κλίση $-1$, άρα
\[
f'(4)=-1.
\]

Επειδή η $f$ έχει συνεχή πρώτη παράγωγο, η $f'$ είναι συνεχής στο
$[1,4]$. Έχουμε
\[
f'(1)=2>0
\qquad\text{και}\qquad
f'(4)=-1<0.
\]
Άρα, από το θεώρημα \eng{Bolzano} για τη συνάρτηση $f'$, υπάρχει
\[
\xi_1\in(1,4)
\]
τέτοιο ώστε
\[
f'(\xi_1)=0.
\]
Επομένως η εφαπτομένη της $C_f$ στο σημείο με τετμημένη $\xi_1$ έχει
κλίση $0$, άρα είναι παράλληλη στον άξονα $x'x$.

\item Η $f$ είναι συνεχής στο $[0,5]$ και παραγωγίσιμη στο $(0,5)$,
αφού έχει συνεχή πρώτη παράγωγο. Επιπλέον, η $C_f$ τέμνει τον άξονα
$x'x$ στα $x=0$ και $x=5$, οπότε
\[
f(0)=0
\qquad\text{και}\qquad
f(5)=0.
\]
Άρα
\[
f(0)=f(5).
\]
Από το θεώρημα \eng{Rolle}, υπάρχει
\[
\eta\in(0,5)
\]
τέτοιο ώστε
\[
{f'(\eta)=0}.
\]
Συνεπώς η εξίσωση $f'(x)=0$ έχει τουλάχιστον μία ρίζα στο $(0,5)$.

\item Θεωρούμε τη συνάρτηση
\[
g(x)=f(x)e^{-x},\quad x\in[0,5].
\]
Η $g$ είναι συνεχής στο $[0,5]$ και παραγωγίσιμη στο $(0,5)$.
Επιπλέον,
\[
g(0)=f(0)e^0=0
\]
και
\[
g(5)=f(5)e^{-5}=0.
\]
Άρα
\[
g(0)=g(5).
\]
Από το θεώρημα \eng{Rolle}, υπάρχει
\[
\xi_2\in(0,5)
\]
τέτοιο ώστε
\[
g'(\xi_2)=0.
\]

Υπολογίζουμε
\begin{align*}
g'(x)
&=\left(f(x)e^{-x}\right)'\\
&=f'(x)e^{-x}-f(x)e^{-x}\\
&=\left(f'(x)-f(x)\right)e^{-x}.
\end{align*}
Άρα
\[
g'(\xi_2)=0
\Leftrightarrow
\left(f'(\xi_2)-f(\xi_2)\right)e^{-\xi_2}=0.
\]
Επειδή
\[
e^{-\xi_2}>0,
\]
προκύπτει
\[
{f'(\xi_2)=f(\xi_2)}.
\]

\item Γνωρίζουμε ότι η $C_f$ τέμνει τον άξονα $x'x$ μόνο στα $x=0$ και
$x=5$. Επομένως η $f$ δεν μηδενίζεται στο $(0,5)$.

Επειδή η $f$ είναι συνεχής και
\[
f(1)=3>0,
\]
η $f$ διατηρεί θετικό πρόσημο σε όλο το $(0,5)$. Άρα
\[
f(x)>0,\quad x\in(0,5).
\]
Συνεπώς
\[
\int_0^5 f(x)\d x>0.
\]
Δηλαδή
\[
{\int_0^5 f(x)\d x>0}.
\]
Γεωμετρικά, στο διάστημα $[0,5]$ η $C_f$ βρίσκεται πάνω από τον άξονα
$x'x$, εκτός από τα άκρα όπου τον τέμνει.
\end{erwthma}
\end{thema}
%=========== ΤΕΛΟΣ ΛΥΣΗΣ - 69ο ΘΕΜΑ Δ ===========
